% CLAUDE WCDC REPORT::WCDC-0001::2026-11-09
\documentclass{article}
\usepackage{amsmath,amssymb,amsthm}
\newtheorem{thm}{Theorem}\newtheorem{prop}[thm]{Proposition}\newtheorem{cor}[thm]{Corollary}
\title{Worst-case derivation cost $\overline W_R(N)$: the bridge dichotomy (WCDC-0001)}
\date{2026-11-09}
\begin{document}\maketitle
\noindent\textbf{Status: P versus NP remains OPEN.} Nothing below proves a lower bound.

\section*{Corpus reviewed}
Root \texttt{AGENTS.md}; \texttt{research/agents/claude/world-ledger} (rows L13, L34--L40); repository searches for ``derivation cost'' found only older LSPV-1 resource-ladder reports, none defining $\overline W_R$. This is the first report on the worst-case-derivation directive. Prior ARC/FLMV and UWI results are untouched and unused here.

\section*{Setup}
$R$ is a Cook--Reckhow proof system (sound, complete for UNSAT, proofs verifiable in time poly$(|\pi|+|F|)$), and $\operatorname{cost}_R(\pi)=|\pi|$ up to polynomial factors. Then
$\overline W_R(N)=\max_{F\in\mathrm{UNSAT},|F|\le N}\min_{\pi}|\pi|$.
\emph{Bridge}$(R)$: $\mathrm{SAT}\in\mathrm P\Rightarrow \overline W_R\le\mathrm{poly}$.

\section*{Results (PROVED/EXACT unless marked)}
\begin{prop}[free strategy $=$ nondeterminism]
The inner $\min$ is the length of the shortest NP-certificate of unsatisfiability in $R$. If $\overline W_R\le\mathrm{poly}$ then $\mathrm{coNP}\subseteq\mathrm{NP}$. Charging for search can only increase cost, so lower bounds on $\overline W_R$ are lower bounds on proof size and nothing more; they are coNP-type statements, not time lower bounds. (KNOWN: Cook--Reckhow 1979.)
\end{prop}
\begin{thm}[bridge dichotomy]
(i) If $\overline W_R$ is unconditionally superpolynomial (resolution, cutting planes, bounded-depth Frege, \dots), then Bridge$(R)\iff\mathrm{SAT}\notin\mathrm P$. So for such $R$ the ``weakest sufficient bridge'' is \emph{equivalent} to the goal, not weaker.\\
(ii) If $R$ is optimal (every proof system is polynomially simulated by $R$), then Bridge$(R)$ holds: $\mathrm P=\mathrm{NP}\Rightarrow\mathrm{NP}=\mathrm{coNP}$ gives a polybounded system $Q$, and $R$ simulates $Q$. But then $\overline W_R\ge\mathrm{superpoly}$ is equivalent to $\mathrm{NP}\ne\mathrm{coNP}$, which is strictly stronger than the target ($\mathrm{NP}\neq\mathrm{coNP}\Rightarrow\mathrm P\ne\mathrm{NP}$ is not known to reverse). Existence of an optimal $R$ is open (Krajíček--Pudlák: follows from $\mathrm{NE}=\mathrm{coNE}$).\\
(iii) For any $R$ with open growth, Bridge$(R)$ requires a black-box transfer: from an \emph{unknown} polynomial-time SAT decider $M$ to poly-size $R$-proofs.
\end{thm}
\begin{proof}
(i) If $\mathrm P=\mathrm{NP}$ then the hypothesis fails to be consistent with superpolynomial $\overline W_R$ only when the conclusion of Bridge fails; so Bridge$(R)$ is the implication $\mathrm P=\mathrm{NP}\Rightarrow\bot$. (ii) as stated. (iii) is a restatement of the quantifier structure.
\end{proof}
\begin{prop}[where the transfer breaks: the UNSAT side]
Let $\mathrm{Th}$ be a theory formalizing polynomial-time computation (e.g.\ PV) and $R_{\mathrm{Th}}$ the system whose proofs are $\mathrm{Th}$-proofs of $\neg\mathrm{Sat}(F)$. If $\mathrm P=\mathrm{NP}$ is witnessed by $M$ with $\mathrm{Th}\vdash\forall F\,(M(F)=0\to\neg\mathrm{Sat}(F))$, then $R_{\mathrm{Th}}$ is polybounded (evaluate $M(F)$ in poly-size, apply the soundness lemma). Hence the weakened bridge
``a $\mathrm{Th}$-provably-correct SAT algorithm exists $\Rightarrow\overline W_{R_{\mathrm{Th}}}\le$ poly'' is true. Levin's universal search is sound-by-construction on the SAT side (it outputs verified assignments), but ``reject $\Rightarrow$ unsatisfiable'' is exactly the part with no known $\mathrm{Th}$-proof. This asymmetry is the whole obstacle. (CANDIDATE reading of the asymmetry; the displayed implication is PROVED.)
\end{prop}
\begin{cor}[what a lower bound would prove]
$\overline W_{R_{\mathrm{Th}}}$ superpolynomial $\Rightarrow$ no $\mathrm{Th}$-provably-correct polynomial SAT decider. This excludes provable algorithms, not all algorithms.
\end{cor}

\section*{Literature clue (this run)}
Pich--Santhanam, \emph{Towards $\mathrm P\ne\mathrm{NP}$ from Extended Frege lower bounds} (arXiv:2312.08163): under circuit-lower-bound hypotheses, an Extended Frege lower bound yields $\mathrm P\ne\mathrm{NP}$. This is a conditional bridge of the (iii)-type for $R=$EF; their hypotheses must be audited (CANDIDATE, not yet read in full). Their companion \emph{Why are proof complexity lower bounds hard?} gives barriers for strong bounds. Classification: useful structural analogy requiring a new transfer theorem (our $\overline W_R$ uses worst-case UNSAT, theirs uses hardness of specified tautologies).

\section*{Red-team notes and next target}
Hypergraph/signed-DAG hard instances must be tested against $R$=EF-level reasoning, not resolution, because by (i) resolution-type bounds are bridge-equivalent to the goal. \textbf{Next target T-BR1:} read arXiv:2312.08163 in full and state exactly which hypothesis replaces the black-box transfer in (iii); decide whether it is implied by, or independent of, SAT$\in$P for a fixed $R$.

\section*{Duplication / change relative to root main}
Nothing here is in root \texttt{main.tex}; no ownership region edited. No claim strengthened.
\end{document}
% END_OF_VERIFIED_REPORT::CLAUDE::WCDC-0001
